\section{Conclusiones}
\label{sec:conclusiones}

\pendiente{Borrador: completar con la latencia medida y con las pruebas de deriva y retorno de la muñeca con la versión actual.}

Se desarrolló un sistema de telerobótica de bajo costo en el que un guante con sensores comanda una mano impresa en 3D por medio de un enlace inalámbrico ESP-NOW. El sistema replica la flexión de los cinco dedos y dos grados de libertad de la muñeca (desplazamiento vertical y rotación), con un firmware sobre FreeRTOS en cada nodo.

Los principales resultados y aprendizajes del proyecto son:

\begin{itemize}
    \item \textbf{Separación de responsabilidades.} El guante concentra la calibración y el cálculo de los ángulos, y la mano sólo valida, suaviza y mueve los servos. Así, cambiar de operador requiere recalibrar únicamente el guante.
    \item \textbf{Seguridad en el firmware.} La verificación por checksum, el modo de seguridad ante la pérdida de paquetes, los límites mecánicos y el limitador de velocidad evitan que un dato erróneo o un corte del enlace dañen la mecánica. La pausa del envío durante la calibración (versión 2.3.0) se apoya en ese mismo modo de seguridad y se verificó en hardware.
    \item \textbf{Proceso iterativo validado en hardware.} Las pruebas de la versión preliminar mostraron que el doble filtrado (en el transmisor y en el receptor) introducía un retardo perceptible y que la implementación del filtro de Kalman producía deriva en la muñeca. Su reemplazo por un filtro complementario sobre ángulos absolutos y por un limitador de velocidad en el receptor resolvió ambos problemas con un algoritmo más simple de ajustar.
    \item \textbf{Respuesta del sistema.} Con la versión 2.3.0, la posición de la mano reproduce la orden del guante en los dedos registrados (pulgar e índice) y en los dos ejes de la muñeca, incluidas las posiciones intermedias; la banda muerta descarta las variaciones de \SI{1}{\degree} de la orden, a costa de un error final de hasta \SI{1}{\degree}.
    \item \textbf{Alimentación.} Las mediciones con la versión 2.3.0 muestran que el receptor consume a \SI{12}{\volt} entre \SIrange{75}{100}{\milli\ampere} en reposo y menos de \SI{0.75}{\ampere} en los picos registrados, muy por debajo de la capacidad de la fuente y del regulador. Separar la alimentación del ESP32 (por USB) de la de los servos evita que un pico de corriente de los motores reinicie el microcontrolador.
\end{itemize}

Entre las limitaciones del sistema actual se encuentran la falta de realimentación de fuerza al operador, la imposibilidad de medir el giro sobre el eje vertical (\textit{yaw}) con el MPU6050, que no tiene magnetómetro, y la dependencia de una calibración inicial de los dedos para cada operador. Además, la latencia de punta a punta se estimó a partir del código y de las hojas de datos, pero todavía no se midió.

\subsection{Trabajo futuro}

Durante el desarrollo se identificaron las siguientes mejoras, que no se implementaron:

\begin{itemize}
    \item \textbf{Medir la latencia real} de punta a punta (sección~\ref{sec:latencia}).
    \item \textbf{Filtro \textit{One Euro}} \cite{casiez2012} en lugar del limitador de velocidad: filtra fuertemente en reposo y casi nada durante los movimientos rápidos.
    \item \textbf{Aumentar la tasa de envío a \SI{100}{\hertz}}. El limitador de velocidad actúa por iteración, por lo que en ese caso debería escalarse con el tiempo real entre iteraciones; de lo contrario la velocidad máxima se duplicaría.
    \item \textbf{PWM de mayor frecuencia en los MG946R}, de \SIrange{100}{200}{\hertz}. El fabricante sólo especifica \SI{50}{\hertz}, por lo que debería probarse vigilando la temperatura de los servos.
    \item \textbf{Robustez de la calibración:} rechazar una calibración en la que la lectura de mano abierta y cerrada coincidan, que hoy produciría una división por cero en la conversión del dedo.
    \item \textbf{Emparejamiento sin reprogramar:} la dirección MAC del receptor está fija en el código del transmisor; si se cambia la placa receptora hay que reprogramar el guante.
    \item \textbf{Realimentación al operador}, por ejemplo háptica, para tareas de manipulación.
\end{itemize}
